\documentclass[aspectratio=169,11pt]{beamer}

\usetheme{Madrid}
\usecolortheme{seahorse}
\setbeamertemplate{navigation symbols}{}
\setbeamertemplate{footline}{\hfill\scriptsize\insertframenumber/\inserttotalframenumber\hspace{1em}\vspace{0.5em}}

\usepackage[utf8]{inputenc}
\usepackage[T1]{fontenc}
\usepackage{amsmath,amssymb}
\usepackage{graphicx}
\usepackage{booktabs}

\title{Coeficiente de Reflexi\'on de un Panel Ligeramente Conductivo}
\subtitle{M\'etodo FDTD vs Matriz de Transferencia}
\author{M\'etodos Computacionales en F\'isica No Lineal}
\institute{M\'aster en F\'isica y Matem\'aticas --- Universidad de Granada}
\date{Curso 2025--2026}

\begin{document}

% ============================================================
\begin{frame}
\titlepage
\end{frame}

% ============================================================
\begin{frame}{Problema y enfoque}

\begin{columns}[T]
\begin{column}{0.55\textwidth}
\textbf{Objetivo:} Calcular $R(f)$ y $T(f)$ de un panel conductivo.

\vspace{0.5em}
\textbf{Métodos:}
\begin{itemize}\setlength\itemsep{0.4em}
  \item TMM analítica
  \item FDTD con fuente temporal y Mur ABC
\end{itemize}

\vspace{0.5em}
\textbf{Conceptos clave:}
\begin{itemize}\setlength\itemsep{0.3em}
  \item $\varepsilon_c = \varepsilon_r - j\sigma/\omega$
  \item $\gamma = j\omega\sqrt{\mu_r\varepsilon_c}$
  \item $\eta = \sqrt{\mu_r/\varepsilon_c}$
  \item Multicapa: $\Phi_{\mathrm{total}} = \Phi_1\Phi_2\cdots\Phi_N$
\end{itemize}

\end{column}

\begin{column}{0.42\textwidth}
\centering
\includegraphics[width=\textwidth]{_figs/esquema.png}

\vspace{0.6em}
{\footnotesize FDTD con Mur ABC, pulso gaussiano hacia la derecha, panel en el centro y observadores antes y después.}
\end{column}
\end{columns}

\end{frame}

% ============================================================
\begin{frame}{Resultados clave}

\begin{columns}[T]
\begin{column}{0.68\textwidth}
\centering
\includegraphics[width=\textwidth]{_figs/comparacion_RT.png}

\vspace{0.3em}
\includegraphics[width=0.48\textwidth]{_figs/barrido_sigma.png}
\hfill
\includegraphics[width=0.48\textwidth]{_figs/multicapa.png}
\end{column}

\begin{column}{0.30\textwidth}
\textbf{Lo más importante:}
\begin{itemize}\setlength\itemsep{0.5em}
  \item FDTD vs TMM: buena concordancia.
  \item Con pérdidas: $|R|^2+|T|^2<1$.
  \item Mayor $\sigma$ aumenta reflexión y reduce transmisión.
  \item Multicapa muestra interferencias más rica.
\end{itemize}
\end{column}
\end{columns}

\end{frame}

% ============================================================
\begin{frame}{Código y validación compacta}

\textbf{Módulos principales:}
\begin{itemize}\setlength\itemsep{0.5em}
  \item \texttt{fdtd1d.py}: \texttt{FDTD1D}, \texttt{set\_panel()}, \texttt{set\_multilayer()}, \texttt{run\_panel\_experiment()}
  \item \texttt{panel\_utils.py}: \texttt{reflection\_transmission()}, \texttt{reflection\_transmission\_stack()}, \texttt{run\_panel\_experiment()}
  \item \texttt{panel\_test.py}: tests unitarios e integración
  \item \texttt{conductive\_panel/}: visualizaciones y comparaciones
\end{itemize}

\vspace{0.8em}
\textbf{Validación breve:}
\begin{itemize}\setlength\itemsep{0.5em}
  \item FDTD simple y multicapa concordante con TMM
  \item Sin pérdidas: energía conservada  aproximada
  \item Con pérdidas: $|R|^2+|T|^2<1$
\end{itemize}

\end{frame}

\end{document}
